

\begin{frame}{The Adjudicative Role: The Main Adversarial Classes}
\hypertarget{src:adjudication}{}
\footnotesize
The procedural \emph{class} fixes who may sue whom, and across the adversarial classes the rule is near-uniform. A \emph{candidate}, \emph{party}, \emph{coalition} or the \textbf{Minist\'erio P\'ublico Eleitoral (MPE)} may file against a \textbf{rival candidate} (on majority slates, also the \emph{vice}); the \textbf{ordinary voter has no standing}. The main classes are:
\smallskip
\begin{itemize}\itemsep3pt
  \item \textbf{AIRC}, registration challenge \ns{(LC 64/90)}: contests a rival's candidacy at registration (ineligibility, \emph{Ficha Limpa}).
  \item \textbf{Representa\c{c}\~ao} \ns{(Lei 9.504)}: campaign-conduct violations, including illicit fundraising, vote-buying (art.~41-A) and misuse of office (art.~73).
  \item \textbf{Propaganda / right of reply} \ns{(art.~58)}: irregular campaign advertising, with the offended side answering.
  \item \textbf{AIJE}, judicial investigation \ns{(LC 64/90 art.~22)}: abuse of economic or political power.
  \item \textbf{AIME / RCED}, mandate and diploma challenge: post-election action to unseat an elected rival.
  \item \textcolor{mygray}{\textbf{Electoral criminal action} \ns{(MPE-led)}: the one \emph{vertical}, institution-driven class.}
\end{itemize}
\smallskip
\scriptsize The arena is therefore \textbf{horizontal}: competitors litigate each other, the \textbf{MPE is the sole institutional actor}, and the voter is absent from the active pole. Judicialization operates here through \textbf{weaponization within the political class}, so the instrument tracks the supply of usable legal instruments across rivals rather than the volume of citizen complaints. \emph{Sources:} Gomes, \emph{Direito Eleitoral} (2024); TSE \emph{Temas Selecionados}.
\end{frame}
